\documentclass[11pt,a4paper]{article}

\usepackage{main-sss}
\usepackage{corrige-sss}

\renewcommand{\matiere}{Physique-Chimie}
\renewcommand{\examen}{DNB}
\renewcommand{\annee}{2025}
\renewcommand{\localisation}{Groupe 1}
\renewcommand{\titresujet}{Ammoniac et engrais}
\renewcommand{\niveau}{Brevet}
\renewcommand{\typedocument}{Corrigé}
\renewcommand{\identifiantsujet}{25GENSCG11}

\begin{document}

\pagination
\titre

\vspace{1em}
\noindent{\large\textbf{Partie 1 : Fabrication de l'ammoniac (14 points)}}

\q
Dans les conditions normales de température et de pression, l'ammoniac est un gaz. La modélisation correcte est la A.

\q
La molécule d'ammoniac NH\textsubscript{3} est composée de 1 atome d'azote (N) et de 3 atomes d'hydrogène (H).

\q
Il s'agit d'une transformation chimique. Les atomes se réorganisent : des molécules disparaissent et une nouvelle molécule apparaît.

\q
Le diazote représente environ 80 \% du volume de l'air (plus précisément 78 \%).

\q
Lors d'une transformation chimique, la masse se conserve : la masse totale des réactifs est égale à la masse totale des produits ($m_\text{diazote} + m_\text{dihydrogène} = m_\text{ammoniac}$).

On en déduit : $m_\text{dihydrogène} = m_\text{ammoniac} - m_\text{diazote} = 1000 - 824 = \SI{176}{kg}$.

Il faut donc \SI{176}{kg} de dihydrogène pour fabriquer une tonne d'ammoniac.

\vspace{1.5em}
\noindent{\large\textbf{Partie 2 : Épandage d'un engrais (11 points)}}

\q

L'action mécanique 1 est modélisée par le schéma B et l'action mécanique 2 est modélisée par le schéma D.

\q
La distance totale parcourue lors des 18 passages est : $d = 18 \times l = 18 \times 250 = \SI{4500}{m} = \SI{4,5}{km}$.

On utilise la formule $v = \dfrac{d}{t}$, soit $t = \dfrac{d}{v}$ :

$t = \frac{\SI{4,5}{km}}{\SI{9}{km/h}} = \SI{0,5}{h}.$

Il faut donc \SI{0,5}{h}, soit \SI{30}{min} pour répandre l'engrais sur la totalité du champ.

\q
Pour traiter les dix champs, la durée nécessaire est : $t_\text{total} = 10 \times \SI{45}{min} = \SI{450}{min} = \SI{7,5}{h}$.

Le forfait A (\SI{5}{h}) est insuffisant car $7,5 > 5$. Le forfait B (\SI{8}{h}) convient car $7,5 < 8$.

L'agriculteur doit donc choisir le forfait B.

\end{document}